\documentclass[12pt,a4paper]{article}

\usepackage{cmap}
\usepackage[T2A]{fontenc}
\usepackage[utf8]{inputenc}
\usepackage[russian]{babel}
\usepackage{amsmath,amssymb}
\usepackage[a4paper,margin=2cm]{geometry}

\newcommand{\R}{\mathbb{R}}
\newcommand{\dd}{\,d}
\newcommand{\Dx}{\Delta x}

\setlength{\parindent}{0pt}
\setlength{\parskip}{0.45em}
\sloppy

\begin{document}

\begin{center}
{\Large\bfseries Билет 4}
\end{center}

\noindent\fbox{%
  \begin{minipage}{\dimexpr\linewidth-2\fboxsep-2\fboxrule\relax}
  \normalfont
Частные производные высших порядков. Независимость смешанной частной
производной от порядка дифференцирования. Дифференциалы высших порядков.
Формула Тейлора для функций нескольких переменных с остаточным членом в формах
Лагранжа и Пеано.
  \end{minipage}%
}

\section*{Краткий план}

Сначала частные производные высших порядков вводятся как последовательное
дифференцирование уже полученных частных производных. Затем формулируется
условие, при котором смешанные производные не зависят от порядка
дифференцирования: непрерывность соответствующих производных в точке. После
этого определяется \(k\)-й дифференциал \(d^k f\) и выводится его координатная
формула. В конце формула Тейлора для функции нескольких переменных получается
из одномерной формулы Тейлора для функции
\(\varphi(t)=f(x^0+t\Dx)\); остаток записывается сначала в форме Лагранжа, а
затем, при непрерывности производных до порядка \(m\), в форме Пеано.

\section*{Определения}

\textbf{Частные производные высших порядков.}
Пусть \(f(x)=f(x_1,\dots,x_n)\) определена в окрестности точки
\(x^0\in\R^n\). Если частная производная
\[
  \frac{\partial f}{\partial x_i}(x)
\]
существует в окрестности \(x^0\), то частная производная этой функции по
\(x_j\) в точке \(x^0\) называется частной производной второго порядка:
\[
  \frac{\partial^2 f}{\partial x_j\,\partial x_i}(x^0)
  =
  \frac{\partial}{\partial x_j}
  \left(\frac{\partial f}{\partial x_i}\right)(x^0).
\]
Частная производная порядка \(k\) определяется индуктивно:
\[
  \frac{\partial^k f}
       {\partial x_{i_k}\cdots\partial x_{i_1}}(x^0)
  =
  \frac{\partial}{\partial x_{i_k}}
  \left(
    \frac{\partial^{k-1}f}
         {\partial x_{i_{k-1}}\cdots\partial x_{i_1}}
  \right)(x^0).
\]
Например, для \(f(x,y)\) производные
\[
  \frac{\partial^2 f}{\partial x\,\partial y},
  \qquad
  \frac{\partial^2 f}{\partial y\,\partial x}
\]
называются смешанными. Без дополнительных условий они могут быть различны.

\textbf{\(k\)-кратная дифференцируемость.}
Функция \(f\) называется \(k\) раз дифференцируемой в точке \(x^0\), если все
частные производные порядка \(k-1\) определены в некоторой окрестности
\(x^0\) и дифференцируемы в точке \(x^0\).

\textbf{Дифференциалы высших порядков.}
Первый дифференциал равен
\[
  df(x^0)=\sum_{i=1}^n
  \frac{\partial f}{\partial x_i}(x^0)\,dx_i.
\]
Дифференциал \(k\)-го порядка определяется по индукции:
\[
  d^k f(x^0)=d(d^{k-1}f)(x^0).
\]
При вычислении \(d^k f\) величины \(dx_1,\dots,dx_n\), входящие в
\(d^{k-1}f\), считаются постоянными, а дифференцирование идет только по
переменным \(x_1,\dots,x_n\).

\textbf{Операторная запись.}
Если \(dx=(dx_1,\dots,dx_n)\), то формально удобно писать
\[
  d^k f(x^0)
  =
  \left(
    dx_1\frac{\partial}{\partial x_1}
    +\dots+
    dx_n\frac{\partial}{\partial x_n}
  \right)^k f(x^0),
\]
где степень означает \(k\)-кратное применение оператора.

\textbf{Многочлен Тейлора.}
Пусть \(f\) имеет производные до порядка \(m\) в окрестности \(x^0\), а
\(\Dx=x-x^0\). Многочлен
\[
  P_m(\Dx)
  =
  f(x^0)+\sum_{k=1}^m\frac{1}{k!}\,d^k f(x^0)
\]
называется многочленом Тейлора порядка \(m\) функции \(f\) в точке \(x^0\).
В этой формуле в дифференциалах полагают \(dx_i=\Delta x_i=x_i-x_i^0\).

\section*{Теоремы}

\textbf{1. Независимость смешанной производной второго порядка.}
Пусть обе смешанные производные
\[
  \frac{\partial^2 f}{\partial x\,\partial y}(x,y),
  \qquad
  \frac{\partial^2 f}{\partial y\,\partial x}(x,y)
\]
определены в окрестности точки \((x^0,y^0)\) и непрерывны в этой точке. Тогда
\[
  \frac{\partial^2 f}{\partial x\,\partial y}(x^0,y^0)
  =
  \frac{\partial^2 f}{\partial y\,\partial x}(x^0,y^0).
\]

\textbf{2. Независимость производных порядка \(k\).}
Если частные производные \(k\)-го порядка функции
\(f(x_1,\dots,x_n)\) определены в окрестности точки \(x^0\) и непрерывны в
\(x^0\), то в этой точке они не зависят от порядка дифференцирования.

\textbf{3. Координатная формула для \(d^k f\).}
Если \(f\) \(k\) раз дифференцируема в точке \(x^0\), то
\[
  d^k f(x^0)
  =
  \sum_{i_1=1}^n\cdots\sum_{i_k=1}^n
  \frac{\partial^k f}
       {\partial x_{i_k}\cdots\partial x_{i_1}}(x^0)
  \,dx_{i_k}\cdots dx_{i_1}.
\]
В частности, при \(n=2\), если производные порядка \(k\) непрерывны в точке,
то
\[
  d^k f(x^0,y^0)
  =
  \sum_{i=0}^k
  \binom{k}{i}
  \frac{\partial^k f}{\partial x^{k-i}\partial y^i}(x^0,y^0)
  \,dx^{k-i}dy^i.
\]

\textbf{4. Формула Тейлора с остатком в форме Лагранжа.}
Пусть \(f(x)=f(x_1,\dots,x_n)\) является \(m+1\) раз дифференцируемой в
некоторой \(\delta\)-окрестности точки \(x^0\). Тогда для любой точки
\(x\in U_\delta(x^0)\)
\[
  f(x)
  =
  f(x^0)
  +\sum_{k=1}^m\frac{1}{k!}\,d^k f(x^0)
  +\frac{1}{(m+1)!}\,d^{m+1}f(x^0+\theta\Dx),
\]
где \(\theta=\theta(x)\in(0,1)\), \(\Dx=x-x^0\), и во всех дифференциалах
\(dx_i=\Delta x_i\).

\textbf{5. Формула Тейлора с остатком в форме Пеано.}
Пусть все частные производные функции \(f\) до порядка \(m\) включительно
существуют и непрерывны в некоторой окрестности точки \(x^0\). Тогда при
\(\Dx=x-x^0\to0\)
\[
  f(x)=P_m(\Dx)+o(|\Dx|^m).
\]

\section*{Доказательства}

\textbf{Независимость смешанных производных второго порядка.}
Рассмотрим сначала функцию двух переменных. Для малого \(t\ne0\) положим
\[
  w(t)=
  f(x^0+t,y^0+t)-f(x^0+t,y^0)-f(x^0,y^0+t)+f(x^0,y^0).
\]
Запишем
\[
  w(t)=\varphi(x^0+t)-\varphi(x^0),
  \qquad
  \varphi(x)=f(x,y^0+t)-f(x,y^0).
\]
По теореме Лагранжа для одной переменной существует
\(\theta_1\in(0,1)\), такое что
\[
  w(t)=
  \left(
    \frac{\partial f}{\partial x}(x^0+\theta_1t,y^0+t)
    -
    \frac{\partial f}{\partial x}(x^0+\theta_1t,y^0)
  \right)t.
\]
Еще раз применяя теорему Лагранжа, теперь к функции
\[
  \psi(y)=\frac{\partial f}{\partial x}(x^0+\theta_1t,y),
\]
получаем для некоторого \(\theta_2\in(0,1)\)
\[
  w(t)=
  \frac{\partial^2 f}{\partial y\,\partial x}
  (x^0+\theta_1t,y^0+\theta_2t)t^2.
\]
По непрерывности смешанной производной
\[
  \lim_{t\to0}\frac{w(t)}{t^2}
  =
  \frac{\partial^2 f}{\partial y\,\partial x}(x^0,y^0).
\]
Если в том же рассуждении поменять роли \(x\) и \(y\), функция \(w(t)\) не
изменится, но получится
\[
  \lim_{t\to0}\frac{w(t)}{t^2}
  =
  \frac{\partial^2 f}{\partial x\,\partial y}(x^0,y^0).
\]
Оба предела равны одному и тому же пределу \(w(t)/t^2\), значит смешанные
производные равны. Для производных порядка \(k\) рассуждают аналогично:
любой порядок дифференцирования получается из другого последовательной
перестановкой соседних дифференцирований, а каждая такая перестановка
допустима по уже доказанному случаю второго порядка и непрерывности
соответствующих производных.

\textbf{Формула для дифференциала \(k\)-го порядка.}
Для \(k=1\) формула совпадает с обычной формулой первого дифференциала.
Для \(k=2\)
\[
  d^2f(x^0)=d(df)(x^0)
  =
  d\left(\sum_{i=1}^n
  \frac{\partial f}{\partial x_i}(x)\,dx_i\right)_{x=x^0}.
\]
Так как \(dx_i\) считаются постоянными,
\[
  d^2 f(x^0)
  =
  \sum_{i=1}^n\sum_{j=1}^n
  \frac{\partial^2 f}{\partial x_j\,\partial x_i}(x^0)\,dx_j\,dx_i.
\]
Повторяя тот же шаг, получаем индукцией
\[
  d^k f(x^0)
  =
  \sum_{i_1=1}^n\cdots\sum_{i_k=1}^n
  \frac{\partial^k f}
       {\partial x_{i_k}\cdots\partial x_{i_1}}(x^0)
  dx_{i_k}\cdots dx_{i_1}.
\]
Отсюда сразу следует операторная запись
\[
  d^k f
  =
  \left(
    dx_1\frac{\partial}{\partial x_1}
    +\dots+
    dx_n\frac{\partial}{\partial x_n}
  \right)^k f.
\]
В случае двух переменных и непрерывных производных порядка \(k\) операторы
\(\partial/\partial x\) и \(\partial/\partial y\) коммутируют, поэтому к
операторной формуле применяется биномиальная формула Ньютона:
\[
  d^k f
  =
  \sum_{i=0}^k
  \binom{k}{i}
  \frac{\partial^k f}{\partial x^{k-i}\partial y^i}
  dx^{k-i}dy^i.
\]

\textbf{Формула Тейлора в форме Лагранжа.}
Зафиксируем \(x\in U_\delta(x^0)\), положим \(\Dx=x-x^0\) и рассмотрим
одномерную функцию
\[
  \varphi(t)=f(x^0+t\Dx),\qquad 0\le t\le1.
\]
Введем оператор
\[
  \widehat A
  =
  \Delta x_1\frac{\partial}{\partial x_1}
  +\dots+
  \Delta x_n\frac{\partial}{\partial x_n}.
\]
По правилу дифференцирования сложной функции
\[
  \varphi'(t)=(\widehat A f)(x^0+t\Dx).
\]
Повторяя дифференцирование, получаем
\[
  \varphi^{(k)}(t)=(\widehat A^k f)(x^0+t\Dx).
\]
По одномерной формуле Тейлора с остатком Лагранжа для \(\varphi\) на
\([0,1]\) существует \(\theta\in(0,1)\), такое что
\[
  \varphi(1)
  =
  \varphi(0)
  +\sum_{k=1}^m\frac{1}{k!}\varphi^{(k)}(0)
  +\frac{1}{(m+1)!}\varphi^{(m+1)}(\theta).
\]
Возвращаясь к \(f\) и используя равенство
\(\widehat A^k f=d^k f\) при \(dx_i=\Delta x_i\), получаем
\[
  f(x)
  =
  f(x^0)
  +\sum_{k=1}^m\frac{1}{k!}d^k f(x^0)
  +\frac{1}{(m+1)!}d^{m+1}f(x^0+\theta\Dx).
\]

\textbf{Формула Тейлора в форме Пеано.}
Применим только что доказанную формулу Лагранжа к порядку \(m-1\):
\[
  f(x)
  =
  f(x^0)
  +\sum_{k=1}^{m-1}\frac{1}{k!}d^k f(x^0)
  +\frac{1}{m!}d^m f(x^0+\theta\Dx).
\]
Достаточно показать, что
\[
  d^m f(x^0+\theta\Dx)-d^m f(x^0)=o(|\Dx|^m).
\]
По координатной формуле
\[
  d^m f
  =
  \sum_{i_1=1}^n\cdots\sum_{i_m=1}^n
  \frac{\partial^m f}
       {\partial x_{i_m}\cdots\partial x_{i_1}}
  dx_{i_m}\cdots dx_{i_1}.
\]
При \(dx_i=\Delta x_i\) каждый моном
\(|dx_{i_m}\cdots dx_{i_1}|\le |\Dx|^m\), а коэффициенты
\[
  \frac{\partial^m f}
       {\partial x_{i_m}\cdots\partial x_{i_1}}(x^0+\theta\Dx)
  -
  \frac{\partial^m f}
       {\partial x_{i_m}\cdots\partial x_{i_1}}(x^0)
\]
стремятся к нулю по непрерывности производных порядка \(m\). Так как сумма
конечная,
\[
  \frac{|d^m f(x^0+\theta\Dx)-d^m f(x^0)|}{|\Dx|^m}\to0.
\]
Следовательно,
\[
  d^m f(x^0+\theta\Dx)
  =
  d^m f(x^0)+o(|\Dx|^m),
\]
и после подстановки получается
\[
  f(x)
  =
  f(x^0)
  +\sum_{k=1}^m\frac{1}{k!}d^k f(x^0)
  +o(|\Dx|^m)
  =
  P_m(\Dx)+o(|\Dx|^m).
\]

\section*{Финальная устная версия}

Частные производные высших порядков получаются последовательным
дифференцированием. Например,
\[
  \frac{\partial^2 f}{\partial x_j\,\partial x_i}(x^0)
  =
  \frac{\partial}{\partial x_j}
  \left(\frac{\partial f}{\partial x_i}\right)(x^0),
\]
а производная порядка \(k\) определяется индуктивно. Смешанные производные в
общем случае могут зависеть от порядка дифференцирования. Основная теорема:
если обе смешанные производные второго порядка определены в окрестности точки
и непрерывны в этой точке, то
\[
  \frac{\partial^2 f}{\partial x\,\partial y}(x^0,y^0)
  =
  \frac{\partial^2 f}{\partial y\,\partial x}(x^0,y^0).
\]
Для производных порядка \(k\) аналогично: если все частные производные
\(k\)-го порядка существуют в окрестности и непрерывны в точке, то результат
не зависит от порядка дифференцирования.

Дифференциалы высших порядков задаются рекурсивно:
\[
  d^k f(x^0)=d(d^{k-1}f)(x^0),
\]
причем при вычислении \(dx_i\) считаются постоянными. Поэтому
\[
  d^k f(x^0)
  =
  \sum_{i_1=1}^n\cdots\sum_{i_k=1}^n
  \frac{\partial^k f}
       {\partial x_{i_k}\cdots\partial x_{i_1}}(x^0)
  dx_{i_k}\cdots dx_{i_1}.
\]
Кратко это записывают так:
\[
  d^k f=
  \left(
    dx_1\frac{\partial}{\partial x_1}
    +\dots+
    dx_n\frac{\partial}{\partial x_n}
  \right)^k f.
\]
Для двух переменных при непрерывности производных порядка \(k\) отсюда по
биному Ньютона получается
\[
  d^k f=
  \sum_{i=0}^k
  \binom{k}{i}
  \frac{\partial^k f}{\partial x^{k-i}\partial y^i}
  dx^{k-i}dy^i.
\]

Многочлен Тейлора порядка \(m\) в точке \(x^0\) равен
\[
  P_m(\Dx)=f(x^0)+\sum_{k=1}^m\frac{1}{k!}d^k f(x^0),
  \qquad \Dx=x-x^0.
\]
Если \(f\) \(m+1\) раз дифференцируема в окрестности \(x^0\), то
\[
  f(x)
  =
  f(x^0)
  +\sum_{k=1}^m\frac{1}{k!}d^k f(x^0)
  +\frac{1}{(m+1)!}d^{m+1}f(x^0+\theta\Dx),
  \qquad \theta\in(0,1),
\]
это остаток в форме Лагранжа. Здесь во всех дифференциалах
\(dx_i=\Delta x_i\). Доказательство сводит задачу к одной переменной:
\(\varphi(t)=f(x^0+t\Dx)\), затем применяется обычная формула Тейлора к
\(\varphi\).

Если частные производные до порядка \(m\) существуют и непрерывны в
окрестности \(x^0\), то остаток можно записать в форме Пеано:
\[
  f(x)=P_m(\Dx)+o(|\Dx|^m),\qquad \Dx\to0.
\]
Это следует из формы Лагранжа порядка \(m-1\): разность
\(d^m f(x^0+\theta\Dx)-d^m f(x^0)\) имеет коэффициенты, стремящиеся к нулю, и
мономы степени \(m\), поэтому она равна \(o(|\Dx|^m)\).

\section*{Источники}

Иванов Г.Е. \emph{Лекции по математическому анализу. Часть 1}:
\texttt{IvGE\_1.pdf}, стр. 191--199.

Основные роли источника: стр. 191--193 -- частные производные высших
порядков и независимость смешанных производных; стр. 193--196 --
дифференциалы высших порядков и операторная запись; стр. 197--199 --
формула Тейлора с остаточными членами в формах Лагранжа и Пеано.

\end{document}
